% ===================================================================
%  LENTELĖ Nr. 12 — Stotis. --- Geležinkeliai. --- Laivai.
%  Traduction 2026 d'apres texto/io, controlee sur texto/fr, texto/ru
%  et texto/cs. Consigne : voir texto/lt/10-lentele-01.tex.
%
%  LE MEME TABLEAU, DANS LA COLONNE TCHEQUE, AVAIT DONNE LA LECON DE
%  LA LIMITE MATERIELLE : « nádraží » avait chasse « banhof », mais
%  « vagón », trop court, avait tenu. J'avais tire de la une regle de
%  syllabes. LE LITUANIEN LA DEFAIT. Il dit « stotis », deux
%  syllabes, qui a chasse « štacija » ; et il dit « vagonas », trois
%  syllabes, qu'il a gardees. La longueur n'y est donc pour rien :
%  l'observation tcheque valait pour le tcheque, non pour les
%  langues.
%
%  CE QUI PARTAGE ICI EST AUTRE CHOSE, ET C'EST NET. TOUT LE
%  MOUVEMENT EST NOMME, TOUT LE MATERIEL EST EMPRUNTE. Les mots
%  lituaniens du chemin de fer sortent tous de cinq verbes :
%
%    « traukinys »   le train, de « traukti » tirer ;
%    « bėgiai »      les rails, de « bėgti » courir ;
%    « garvežys »    la locomotive, le porte-vapeur, de « vežti » ;
%    « stotis »      la gare, de « stoti » s'arreter ;
%    « kūrikas »     le chauffeur, de « kūrenti » chauffer ;
%    « iešmininkas » l'aiguilleur, de « iešmas » la broche.
%
%  Restent etrangers « vagonas », « peronas », « bilietas »,
%  « tunelis », « semaforas », « konduktorius », « bufetas ». Et
%  « geležinkelis », le chemin-de-fer, est un calque exact — la
%  troisieme voie tcheque, une fois de plus.
%
%  C'EST LA MEME PARTITION QU'AU CHANTIER DU TABLEAU 5, et elle se
%  verifie donc sur deux techniques arrivees tard. La langue nomme ce
%  qu'elle voit FAIRE et emprunte ce qu'elle voit ARRIVER TOUT FAIT.
%  Un rail court, une machine tire, un homme chauffe : cela se dit en
%  lituanien. Un wagon, un quai, un billet : cela vient avec la
%  chose.
%
%  ET LA NOTE DU FEUILLET 85 EST DE 1926, non de 1914 : elle avertit
%  que le voyage est decrit « segun la pre-milita frontiero », selon
%  la frontiere d'avant-guerre. On la rend telle quelle, comme les
%  treize colonnes precedentes.
% ===================================================================

%%K t12-apar-1 apar
\VUsaut{4.0mm}
\VUtitre{15.0pt}{60}{LENTELĖ Nr. 12}
\VUsaut{2.4mm}
\VUpk{11.4pt}{40}{Stotis. --- Geležinkeliai. --- Laivai.}
\VUsaut{3.4mm}
\textsc{Pastaba.} --- \textit{Kiekvienoje tautoje vartojami tvarkaraščiai yra
skirtingi, todėl pavyzdžiu imame} \VUgras{prancūziškos lentelės} \textit{valandas.}

%%K t12-01-1 p
1. --- Ką tik keliavau po Vokietiją. Prieš išvykdamas nusipirkau \VUgras{tvarkaraščio
knygelę} \textsuperscript{(15)}, kad žinočiau traukinių išvykimo valandą. Įsitikinęs,
kad patogiausias traukinys išvyksta iš Paryžiaus devintą valandą vakaro ir atvyksta į
Strasbūrą pirmą valandą trylika minučių, pasiruošiau kelionę, ir kitą dieną liepiau
nuvežti mane į stotį.

%%K t12-02-1 p
2. --- Kai įėjau į didžiąją išvykimo salę, už seno kaimo \VUgras{kunigo}
\textsuperscript{(2)}, kuris rankoje nešėsi savo \VUgras{kelionmaišį}
\textsuperscript{(3)}, jau buvo atvykę daug \VUgras{keleivių} \textsuperscript{(1)}.

%%K t12-02-2 p
Kadangi norėjau pasiųsti \VUgras{depešą (telegramą)} \textsuperscript{(5)}, paprašiau
\VUgras{kontrolieriaus} \textsuperscript{(6)}, kad nurodytų man \VUgras{telegrafo
kontorą} \textsuperscript{(4)}.

%%K t12-03-1 p
3. --- Prieš eidamas į \VUgras{laukiamąją} \textsuperscript{(7)}, nuėjau apsirūpinti
\VUgras{bilietu} \textsuperscript{(9)} su grįžimu.

%%K t12-04-1 p
4. --- Neilgai laukiau prie \VUgras{langelio} \textsuperscript{(8)}, nes ten buvo
nedaug žmonių. \VUgras{Kareivis} \textsuperscript{(10)}, kuris išvyko atostogų,
malonėjo užleisti man savo eilę, tad turėjau pakankamai laiko paklausti keleto žinių
\VUgras{stoties viršininko} \textsuperscript{(13)}, kuris kalbėjosi su priežiūros
\VUgras{komisaru} \textsuperscript{(12)} ir su budinčiu \VUgras{žandaru}
\textsuperscript{(11)}.

%%K t12-tit-1 sub
\VUpk{10.2pt}{40}{Bagažo įregistravimas.}
\VUsaut{3.0mm}

%%K t12-05-1 p
5. --- Nusipirkęs laikraštį \VUgras{knygų kioske} \textsuperscript{(14)}, liepiau
pasverti savo \VUgras{bagažą} \textsuperscript{(17)}, kurį \VUgras{nešikas}
\textsuperscript{(16)} ką tik atvežė \VUgras{bagažo vežimėliu} \textsuperscript{(19)};
\VUgras{tarnautojas} \textsuperscript{(22)}, kuriam pavestos \VUgras{svarstyklės}
\textsuperscript{(21)}, pranešė man, kad mano lagaminas sveria trisdešimt penkis
kilogramus; tai davė penkių kilogramų \VUgras{antsvorio,} už kurį turiu papildomai
sumokėti prie \VUgras{bagažo langelio} \textsuperscript{(23)}.

%%K t12-06-1 p
6. --- Kadangi buvau per anksti, turėjau laisvo laiko stebėti keleivių judėjimą
stotyje. Vieni padėjo ar atsiėmė savo nesvarbų bagažą \VUgras{saugykloje}
\textsuperscript{(25)}; kiti žiūrėjo į \VUgras{tvarkaraščius} \textsuperscript{(26)}.
Atvykę keleiviai skubėjo prie \VUgras{išėjimo} \textsuperscript{(32)} su savo
\VUgras{lagaminu} \textsuperscript{(24)} rankoje. Kai kurie laukė \VUgras{bagažo
išdavimo vietoje} \textsuperscript{(27)} ir rodė savo \VUgras{kvitą}
\textsuperscript{(29)}, kad atsiimtų savo lagaminus, \VUgras{dėžes}
\textsuperscript{(28)} ar \VUgras{krepšius} \textsuperscript{(30)}.

%%K t12-07-1 p
7. --- \VUgras{Akcizo pareigūnai} \textsuperscript{(33)} rūpestingai tikrino ryšulius,
kad įsitikintų, ar nėra ko deklaruoti.

%%K t12-08-1 p
8. --- Kai kurie keleiviai nuėjo į \VUgras{bufetą} \textsuperscript{(34)}, kur
\VUgras{padavėjas} \textsuperscript{(35)} stengėsi juos aptarnauti. Jiems reikia
skubėti, nes \VUgras{traukinys} \textsuperscript{(36)}, kuris yra ekspresas, neilgai
stovės stotyje.

%%K t12-09-1 p
9. --- Paskui nuėjau į išvykimo peroną ir ieškojau tiesioginio \VUgras{vagono}
\textsuperscript{(37)}, kad nebūčiau priverstas keisti vagono kelionės metu. Įlipau į
koridorinį vagoną, kuriame yra \VUgras{skyrius} \textsuperscript{(38)} rūkantiesiems
ir skyrius, skirtas „vienišoms ponioms“.

%%K t12-tit-2 sub
\VUpk{10.2pt}{40}{Išvykimas naktį.}
\VUsaut{3.0mm}

%%K t12-10-1 p
10. --- Užlipęs ant \VUgras{pakopos} \textsuperscript{(40)}, uždaręs \VUgras{dureles}
\textsuperscript{(39)}, susukęs \VUgras{užuolaidėlę} \textsuperscript{(41)} ir padėjęs
savo mažą bagažą ant \VUgras{tinklelio} \textsuperscript{(43)}, atsisėdau ant
\VUgras{suolelio} \textsuperscript{(42)}. Matydamas \VUgras{žibintininką}
\textsuperscript{(44)}, uždegantį lempas, pagalvojau, kad turėsiu praleisti naktį
vagone, ir pasišaukiau tarnautoją, kuriam pavesta nuomoti \VUgras{pagalvėles}
\textsuperscript{(45)}, kad paprašyčiau vienos.

%%K t12-11-1 p
11. --- Traukinio konduktorius atėjo tikrinti bilietų. Paklausiau jo, kodėl dar
neišvykome, nes jau tikroji valanda. Jis man atsakė, kad \VUgras{traukinio
viršininkas} \textsuperscript{(46)} užsiėmęs, liepdamas sudėti dviračius į
\VUgras{furgoną} \textsuperscript{(47)}. Be to, dar nebuvo pakeisti visi kojų
šildytuvai \textsuperscript{(48)}.

%%K t12-12-1 p
12. --- Bet netrukus turėjome išvykti, nes \VUgras{kūrikas} \textsuperscript{(50)}
ant savo \VUgras{tenderio} \textsuperscript{(49)} paėmė anglių, kad pakurstytų
\VUgras{garvežio} \textsuperscript{(54)} ugnį, o \VUgras{mašinistas}
\textsuperscript{(51)} telaukė \VUgras{stoties viršininko padėjėjo}
\textsuperscript{(52)} ženklo, kuris buvo ant \VUgras{perono} \textsuperscript{(53)}.

%%K t12-13-1 p
13. --- Garvežio \VUgras{katilas} \textsuperscript{(56)} duslai gaudė;
\VUgras{kaminas} \textsuperscript{(55)} vėmė dūmų srautus, sumaišytus su
\VUgras{garais} \textsuperscript{(60)}. Skardus \VUgras{švilpukas}
\textsuperscript{(59)} nuaidėjo, ir \VUgras{stūmokliai} \textsuperscript{(58)}
pajudėjo, stumdami sunkios mašinos \VUgras{ratus} \textsuperscript{(57)}.

%%K t12-tit-3 sub
\VUsaut{3.0mm}
\VUpk{10.2pt}{40}{Išvykstame!}
\VUsaut{3.0mm}

%%K t12-14-1 p
14. --- Traukinio, bėgančio \VUgras{bėgiais} \textsuperscript{(64)}, greitis didėjo;
\VUgras{peronai} \textsuperscript{(65)} dingo; pravažiavome pro \VUgras{išvietes}
\textsuperscript{(66)}, taip pat pro vagonų sąstatą, pastatytą ant kito
\VUgras{kelio} \textsuperscript{(63)}: \VUgras{miegamuosius vagonus}
\textsuperscript{(67)}, \VUgras{restorano vagoną} \textsuperscript{(68)},
\VUgras{pašto vagoną} \textsuperscript{(69)}.

%%K t12-noto-1 noto
\VUnotes{143.2mm}{%
\textit{(*) Pastaba.} --- \textit{Savaime suprantama, kelionė aprašoma pagal prieškarinę
sieną.}%
}

%%K t12-15-1 p
15. --- Paskui pro mūsų akis praslinko \VUgras{prekių stotis} \textsuperscript{(71)}.
Po angarais tarnautojai krovė \VUgras{prekes} \textsuperscript{(72)}, siunčiamas
negreituoju traukiniu. \VUgras{Brigadininkai} \textsuperscript{(76)} prie
\VUgras{vandens rezervuaro} \textsuperscript{(73)} manevravo \VUgras{gyvulinius
vagonus} \textsuperscript{(75)} ir \VUgras{platforminius vagonus}
\textsuperscript{(79)}, kad sudarytų \VUgras{prekinį traukinį} \textsuperscript{(80)}.

%%K t12-16-1 p
16. --- Pervažiavome paskutinį \VUgras{išsišakojimą} \textsuperscript{(78)}, prie
kurio stovėjo iešmininkas; pravažiavome \VUgras{signalinį diską}
\textsuperscript{(86)}, ir stotis toli dingo.

%%K t12-17-1 p
17. --- Netrukus pervažiavome \VUgras{geležinį tiltą} \textsuperscript{(74)}, kuris
kerta \VUgras{upę} \textsuperscript{(81)}. Pervažiavę tą tiltą, važiavome palei upę,
kuria galingi \VUgras{vilkikai} \textsuperscript{(82)} vilko sunkias \VUgras{baržas}
\textsuperscript{(83)}. Dar matėme \VUgras{keleivinį laivą} \textsuperscript{(84)},
kai jis plaukė prie \VUgras{pontono} \textsuperscript{(85)}.

%%K t12-18-1 p
18. --- Greitai pravažiavę kelias stotis, pervažiavome keletą \VUgras{tunelių}
\textsuperscript{(87)} ir palikome už savęs nesuskaičiuojamus \VUgras{užtvarų sargų}
\VUgras{namelius} \textsuperscript{(88)}. Supamas traukinio kratymo, giliai užmigau.

%%K t12-tit-4 sub
\VUsaut{3.0mm}
\VUpk{10.2pt}{40}{Deutsch-Avricourt stotis.}
\VUsaut{3.0mm}

%%K t12-19-1 p
19. --- Mane staiga pažadino tarnautojo balsas, šaukiantis: „Deutsch-Avricourt!
\textsuperscript{(*)}. Visi lipkite!“ Įvyko muitinės patikra ir visos varginančios
pasienio formalybės. Turėjau suderinti savo kišeninį laikrodį su stoties
\VUgras{laikrodžiu} \textsuperscript{(89)}, kuris skubėjo penkiasdešimt penkiomis
minutėmis; vos pažadintas, sunkiai skyriau \VUgras{didžiąją rodyklę}
\textsuperscript{(90)} nuo \VUgras{mažosios} \textsuperscript{(91)}; pats
\VUgras{ciferblatas} \textsuperscript{(92)} buvo visai neaiškus dėl rytmečio rūko.

%%K t12-20-1 p
20. --- Kol muitininkai darė savo patikrą, žiovaudamas skaičiau \VUgras{afišas}
\textsuperscript{(93)}, kurios puošia stoties sienas. Papusryčiavau, kad praleisčiau
laiką, pažiūrėjau į Vokietijos \VUgras{tinklo} \textsuperscript{(94)} žemėlapį, kad
pamatyčiau, koks nuotolis mane dar skiria nuo Strasbūro, ir grįžau į savo vagoną,
sustiprintas vilties netrukus pasiekti savo kelionės tikslą.

\VUsaut{6.0mm}
\VUfilet{22mm}
